\subsection{向量空间中的线性方程}

命题 12. 设 E, F 是域 K 上的两个向量空间，且 u:E → F 是一个线性映射。对于线性方程

\[
u (x) = y _ {0}\tag{20}
\]

至少有一个解 \(x \in \mathbf{E}\) ，必要且充分的条件是 \(y_0\) 与转置映射 \(^t u\) 的核正交.

\(u(\mathbf{E})\) 在 \(\mathbf{F}^*\) 中的正交补是 \({}^t u^{-1}(0)\) （§ 2，第 5 号，命题 8 的推论），而 \({}^t u^{-1}(0)\) 在 \(\mathbf{F}\) 中的正交补因此为 \(u(\mathbf{E})\) （第 5 号，定理 7 (ii)）。

我们将得到关于标量线性方程组

\[
\langle x, x _ {\iota} ^ {*} \rangle = \eta_ {\iota} \quad (\iota \in I)\tag{21}
\]

的一个更便捷的判别准则，其中未知量 x 在域 K 上的向量空间 E 中取值， \(x_{\iota}^{*}\) 是 E 上的线性型，而右端项 \(\eta_{\iota}\) 是 K 中的元素。如果我们考虑 E 的一组基

\((a_{\lambda})_{\lambda\in L}\) ，则方程组 (21) 等价于方程组

\[
\sum_ {\lambda \in L} \xi_ {\lambda} \left\langle a _ {\lambda}, x _ {\iota} ^ {*} \right\rangle = \eta_ {\iota} \quad (\iota \in I)\tag{22}
\]

其中 \(x = \sum_{\lambda \in L} \xi_{\lambda} a_{\lambda}\) ，而 (22) 的解必然是族 \((\xi_{\lambda})\) ，其元素属于 \(K\) 且具有有限支撑。

定义 4. \(\mathbf{E}^*\) 中由族 \((x_{\iota}^{*})\) 生成的子空间的维数称为系统 (21) 的秩。

命题 13. 系统 (21) 具有有限秩 \(r\) ，必要且充分条件是线性映射 \(u: x \mapsto (\langle x, x_{\iota}^* \rangle)\) 把 \(\mathbf{E}\) 映入 \(\mathbf{K}_s^I\) 且其秩为 \(r\) .

若 \(F'\) 是 \(E^{*}\) 中由诸 \(x_{\iota}^{*}\) 生成的子空间，则 u 的核是 \(F'\) 在 E 中的正交补 F；如果 \(F'\) 的维数为 r，则 F 的余维数为 r，反之亦然（第 5 号，定理 7），且 \(\operatorname{rg}(u) = \operatorname{codim}_{\mathbf{E}} \mathbf{F}\) （第 4 号，公式 (11)）。

定理 8. 设

\[
\langle x, x _ {\iota} ^ {*} \rangle = \eta_ {\iota} \quad (\iota \in I)\tag{21}
\]

设 E 是域 K 上的向量空间，考虑 E 上的一个标量线性方程组。要使该方程组至少有一个解，必须满足：对任意具有有限支撑的标量族 \((\rho_{\iota})\) ，只要 \(\sum_{\iota} x_{\iota}^{*}\rho_{\iota}=0\) ，就有 \(\sum_{\iota}\eta_{\iota}\rho_{\iota}=0\) 。如果系统 (21) 的秩是有限的，那么这个条件也是充分的。

该条件显然是必要的。它表明，如果 \(F'\) 是 \(E^{*}\) 中由族 \((x_{i}^{*})\) 生成的子空间，则存在一个线性映射 \(f:F'\to K_{d}\) ，使得 \(f(x_{\iota}^{*})=\eta_{\iota}\) 对所有 \(\iota\in I\) 成立。若 \(F'\) 是有限维的，维数为 r，则 \(F'\) 是 E 的余维数为 r 的子空间 F 的正交补（第 5 号，定理 7），并且 \(F'\) 等同于 E/F 的对偶空间（§ 2，第 6 号，命题 9 的推论）；因此 f 是双对偶空间 \((\mathrm{E}/\mathrm{F})^{**}\) 中的一个元素。由于 E/F 是有限维的，存在唯一一个元素 \(y\in E/F\) ，使得 \(f(x^{*})=\langle y,x^{*}\rangle\) 对所有 \(x^{*}\in F'\) 成立（第 5 号，定理 6）。于是 (21) 的解即为那些在 E/F 中的典范像为 y 的 \(x\in E\) 。

注记。当系统 (21) 的秩为无穷时，定理 8 的条件不再充分。例如，假设 \(x_{\iota}^{*}\) 是空间 \(\mathbf{E} = \mathbf{K}_{s}^{(\mathrm{I})}\) 上的坐标线性型，其中 I 是无限的（§ 2，第 6 号）；由于 \(x_{\iota}^{*}\) 线性无关，定理 8 的条件对每个族 ( \(\eta_{\iota}\) ) 都成立，但系统 (21) 仅当族 ( \(\eta_{\iota}\) ) 具有有限支撑时才有解。

系统 (21) 若只含有限个方程，则其秩总是有限的，且此时秩至多等于方程的个数（第 2 号，命题 1）。类似地，若 E 为有限维，维数为 n（对于系统 (22) 而言，这对应于未知量仅有有限个 n 的情形），则其对偶 \(E^{*}\) 的维数为 n，因此系统 (21) 的秩至多为 n（第 3 号，命题 4 的推论 4）。由此我们推出：

推论 1. 向量空间中由有限个方程组成的标量线性方程组，若各方程的左端是线性无关的线性型，则恒有解。

推论 2. 对于系数在域 K 中、含 n 个未知量的齐次线性方程组 (22)，要使它存在由 K 中元素组成的非平凡解，必要且充分条件是其秩 \textless n.

如果方程的个数是有限的并且 \textless n ，那么情况总是如此。

推论 3. 对于系数和右端项均属于域 K 的线性方程组 (22)，若其包含 n 个方程和 n 个未知数，则它有且仅有一个由 K 中元素构成的解，当且仅当相应的齐次方程组没有非平凡解（或者，等价地说，该方程组的左端是线性无关的线性型）。

